\documentclass[11pt,a4paper]{article}
\usepackage{amsmath,amssymb,amsthm}
\usepackage[margin=2.5cm]{geometry}
\usepackage{hyperref}

\newtheorem{definition}{Definition}[section]
\newtheorem{theorem}[definition]{Theorem}
\newtheorem{proposition}[definition]{Proposition}
\newtheorem{lemma}[definition]{Lemma}
\newtheorem{corollary}[definition]{Corollary}
\newtheorem{conjecture}[definition]{Conjecture}
\newtheorem{remark}[definition]{Remark}

\title{Hyperseed Formalization:\\
  Linguistic Universals as Shadows of Cognitive Structure, Part 1}
\author{ProtomegaTron (automated formalization)}
\date{2026-09-18}

\begin{document}
\maketitle

\begin{abstract}
We formalize the intellectual structure of Ben Goertzel's Substack article
``Linguistic Universals as Shadows of Cognitive Structure, Part 1''
(2026-05-20) within the Hyperseed ontology. The article establishes that
linguistic universals are structural fingerprints of cognitive architecture,
grounded in Verkerk et al.'s phylogenetic analysis. Five conclusions about
grammar space (closure systems, mediation requirement, coherence laws,
graded universalhood, context-indexing) each map onto features of
causal-coding neural architecture (HBCML/ColBaC) and onto Hyperseed
structures: fiber closure, fiber product, sparse fiber routing, fiber
depth, and sheaf structure.
\end{abstract}

\tableofcontents

%% =================================================================
\section{Empirical grounding: the phylogenetic filter}
\label{sec:empirical}
%% =================================================================

\begin{proposition}[Verkerk et al.\ survival rates --- claims 1--6]
\label{prop:empirical}
The article's empirical foundation is Verkerk et al.'s Bayesian
phylogenetic analysis of Greenberg's universals, controlling for
genealogical and areal confounds. The survival rates across
categories are:

\begin{center}
\begin{tabular}{lcc}
\textbf{Category} & \textbf{Survive} & \textbf{Rate} \\
\hline
Hierarchical universals & 24/30 & 80\% \\
Narrow word-order & 24/65 & 37\% \\
Broad word-order & 8/72 & 11\% \\
Miscellaneous & 4/24 & 17\% \\
\end{tabular}
\end{center}

The hierarchical class dominates. Among survivors, hierarchical and
narrow word-order universals show directed drift toward harmonic
states under phylogenetic modeling; broad word-order claims show no
such asymmetry. The article takes this pattern as ``too extremely
vivid to be noise'' and asks what mathematical structure would
naturally generate exactly this profile.
\end{proposition}

%% =================================================================
\section{C1: Hierarchical universals as closure systems (fiber closure)}
\label{sec:closure}
%% =================================================================

\begin{theorem}[Closure as attractor --- claims 7--10, 29]
\label{thm:closure}
Hierarchical universals share a cumulative structure: rarer or more
marked features only appear together with the more common ones
beneath them. Person hierarchy: if a language has verbal agreement
with third person, it also has agreement with first and second.
The Keenan--Comrie accessibility hierarchy works the same way.

The mathematical structure is a \emph{closure system}: well-formed
paradigms are downward-closed regions of a partial order:
\[
  f \in \text{Paradigm} \;\wedge\; f' \leq f \implies
  f' \in \text{Paradigm}
\]

The closure region is an \emph{attractor in grammar space}: Verkerk
et al.'s finding of directed drift toward harmonic states is the
attractor's signature. Languages drift toward closure-satisfying
states more often than away.

In Hyperseed terms, hierarchical universals are \emph{fiber closure}:
downward-closed regions on the cognitive fiber that are attractors
under diachronic dynamics. The closure is a fiber property ---
it persists across all sections (languages) because it belongs to
the fiber (cognitive substrate), not to any particular section
(individual language). The diachronic drift is the fiber's
attractor dynamics manifesting through the sections.
\end{theorem}

%% =================================================================
\section{C2: Cross-module universals require mediation (fiber product)}
\label{sec:mediation}
%% =================================================================

\begin{theorem}[Product blocks cross-factor implications --- claims 11--13, 30]
\label{thm:mediation}
Broad universals link features across grammatical modules (syntax
with morphology, phonology with syntax). These are the weakest
empirical category. The mathematical reason: if two modules are
genuinely independent --- operations in one can be permuted with
operations in the other without changing observable output --- the
global grammar factors as their product:
\[
  G \cong G_1 \times G_2 \quad\implies\quad
  \nexists \text{ nontrivial } G_1\text{-}G_2 \text{ implication}
\]

Any genuine broad universal must violate at least one independence
assumption, typically through a \emph{mediator} --- case marking,
argument-structure resolution, agreement morphology --- that shares
resources with both modules.

In Hyperseed terms, cross-module independence is \emph{fiber product
structure}: independent cognitive modules are independent fibers
whose product is trivial. Genuine broad universals require
\emph{non-trivial transition functions} between fibers --- mediators
that break the product structure. The fragility of broad universals
is predicted by the prevalence of approximate product structure in
cognition.
\end{theorem}

%% =================================================================
\section{C3: Word-order coherence as sparse fiber routing}
\label{sec:wordorder}
%% =================================================================

\begin{proposition}[Sparse signed energy --- claims 14--16, 31]
\label{prop:wordorder}
Word order is not a single global head-direction parameter, nor
arbitrary. It's a network of locally coupled choices: SOV languages
tend toward postpositions, genitive-noun, relative-noun,
auxiliary-after-verb. SVO and VSO lean the opposite way. Mixed
types (Persian: SOV but prepositional) pay a coherence cost but
are clearly possible.

The model: a sparse signed-energy network where each pair of
word-order choices has a preferred relationship. Real languages
cluster in low-energy configurations; language change biases
drift toward lower-energy states.

In Hyperseed terms, word-order coherence is \emph{sparse fiber
routing}: few high-traffic pathways between fibers rather than
dense coupling everywhere. The energy landscape is the routing
structure --- low-energy configurations correspond to coherent
routing patterns. This is anti-scalar-collapse applied to
inter-module coupling: the word-order structure is a small number
of identifiable pairwise preferences, not a collapsed global
parameter.
\end{proposition}

%% =================================================================
\section{C4: Graded universalhood as fiber depth}
\label{sec:graded}
%% =================================================================

\begin{definition}[Stability regimes --- claims 17--18, 32]
\label{def:graded}
Universalhood is graded, not binary. Each pattern has a stability
coefficient measuring how strongly it resists perturbation across
contexts and languages. Three regimes:

\begin{enumerate}
\item \textbf{Rigid kernel invariants (highest stability).}
  Hierarchical universals. Protected by the deepest cognitive
  substrate. Extremely slow to change.
\item \textbf{Consolidated preferences (moderate stability).}
  Narrow word-order universals. Stable enough to show directional
  drift but not as deeply protected as hierarchies.
\item \textbf{Context-local residue (lowest stability).} Where
  most actual variation lives. High plasticity, fast adaptation.
\end{enumerate}

In Hyperseed terms, this is \emph{fiber depth}: the protected
kernel is the deep fiber layer (high stability, slow learning);
adaptable shells are shallow layers (low stability, fast learning).
The same graded-fiber-depth structure that appears in the four-layer
security architecture (AGI Catastrophe Part 2) and the cognitive
training recipe (Part 3), applied here to the empirical typological
record.
\end{definition}

%% =================================================================
\section{C5: Context-indexed universals as sheaf structure}
\label{sec:context}
%% =================================================================

\begin{definition}[Local closure operators --- claims 19--20, 33]
\label{def:context}
Some universals hold only within a context class (e.g., ``among
SOV languages with case marking...''). These are local closure
operators, not global ones.

The local closure operators are glued by transition maps across
contexts, forming a sheaf-like structure: each context has its own
closure system, and the transition maps between contexts preserve
the closure structure on overlapping regions.

In Hyperseed terms, this is \emph{context-indexed fiber structure}:
local sections (context-specific closure operators) glued by
transition maps into a sheaf. The sheaf condition --- that local
sections agree on overlapping contexts --- is the consistency
condition on the fiber bundle's transition functions.

This is the most sophisticated structural feature: it says that
cognitive architecture is not globally uniform but locally
structured with compatible transition rules. The ``universals''
are properties of the sheaf (global sections or near-global
sections), not properties of any single local chart.
\end{definition}

%% =================================================================
\section{Causal coding as fiber commutativity}
\label{sec:causal}
%% =================================================================

\begin{theorem}[HBCML/ColBaC correspondence --- claims 21--28, 34]
\label{thm:causal}
The article establishes a point-by-point correspondence between
the five linguistic conclusions and the structural features of
causal-coding neural architecture (HBCML/ColBaC):

\begin{center}
\begin{tabular}{lll}
\textbf{Linguistic} & \textbf{Cognitive} & \textbf{Hyperseed} \\
\hline
C1: Closure & Hard kernel & Fiber closure \\
C2: Mediation & Factorization & Fiber product \\
C3: Coherence & Sparse routing & Sparse fiber routing \\
C4: Graded & Kernel-shell & Fiber depth \\
C5: Context-indexed & Context-gating & Sheaf structure \\
\end{tabular}
\end{center}

The central result: the \emph{causal coding theorem}. When a
gradient-based learning system maintains causal factorization,
gradient updates triggered by different contexts approximately
commute:
\[
  \nabla_{A} \circ \nabla_{B} \approx \nabla_{B} \circ \nabla_{A}
  \qquad \text{when } \text{supp}(A) \cap \text{supp}(B) = \emptyset
\]

Catastrophic forgetting is bounded by module overlap --- and the
bound goes to zero when supports are disjoint.

In Hyperseed terms, this is the \emph{fiber commutativity theorem}:
updates to independent fibers don't interfere. Catastrophic
forgetting is a \emph{fiber defect} --- a failure of the
decomposition to be genuine. The causal coding theorem provides the
structural guarantee that the fiber decomposition is well-defined:
independent fibers are genuinely independent, and learning in one
doesn't corrupt another.

The kernel-shell architecture (hard kernel defended across contexts,
adaptable shells absorbing local adaptation, outermost shell
carrying disposable residue) is fiber depth implemented in neural
hardware. The hard kernel is the deep fiber; the shells are
progressively shallower layers; the outermost shell is the shallowest.
\end{theorem}

%% =================================================================
\section{Cross-references}
%% =================================================================

\begin{remark}[Connections to other formalizations]
\begin{itemize}
\item \textbf{Linguistic Universals Part 2} (2026-05-30): companion
  article developing the categorical formalism (quantale weakness,
  pushforward theorems, the five correspondences sharpened).
\item \textbf{Linguistic Universals Part 3} (2026-06-08): companion
  article drawing out the AGI architecture implications and concrete
  Transformer training recipe.
\item \textbf{Avoiding AGI Catastrophe Part 2} (2026-06-10): graded
  fiber depth. The four-layer security architecture has the same
  structural pattern as the three-stability-regime linguistic pattern.
\item \textbf{Seeding RSI} (2026-07-28): Hyperon architecture. The
  HBCML/ColBaC architecture is a complement to Hyperon's
  neural-symbolic approach.
\item \textbf{The Orchard Bug} (2026-06-05): cocycle defects. The
  mediation requirement (C2) is the positive version of the Orchard
  bug lesson: without explicit mediators, cross-module composition
  fails silently.
\item \textbf{$(f,c)$-lossy} (LTM 5a4d150b): anti-scalar-collapse.
  Sparse fiber routing (C3) and the graded universalhood (C4) are
  anti-scalar-collapse applied to linguistic structure.
\end{itemize}
\end{remark}

\end{document}
